\clearpage\fancyhead[L]{\small COMET.}\SourceLeaf{346}
\begin{center}\textbf{COMET.}\par\vspace{4pt}\hrule\vspace{4pt}\fontsize{10}{12}\selectfont Monatliche Correspondenz. Band XXIV. S. 507\,..\,517. 1811 November.\par\vspace{4pt}\hrule\vspace{4pt}\end{center}

\begin{flushright}Göttingen den 15. November.\end{flushright}
\par\indent Nach dem was \textsc{Pingré} von europäischen und chinesischen Beobachtungen des Cometen von 1301 anführt, \emph{kann dieser mit dem gegenwärtigen nicht identisch sein.} Denn es ist unmöglich, dass letzterer am 16. Sept. unter einer Länge von 110$^\circ$ und eben so unmöglich, dass er am 30. Sept. unter einer Länge von 231$^\circ$ und \emph{nur} 26$^\circ$ nördl. Breite erschienen wäre.


\par\indent Theils ungünstiges Wetter, theils andere Geschäfte, haben mir seit meinem letzten Brief nur eine Beobachtung des Cometen erlaubt, die ich Ihnen hier mittheile:

\begin{center}
\fontsize{10}{12}\selectfont\setlength{\tabcolsep}{3pt}\renewcommand{\arraystretch}{1.15}\begin{tabular}{l|c|c}
 & AR. & Decl. bor. \\ \hline
1811 Oct. 20 $10^{\mathrm u}$ $9^{\mathrm m}$ $12^{\mathrm s}$ m. Z. & $247^\circ\,46'\,39''$ & $41^\circ\,20'\,15''$ \\
\end{tabular}
\end{center}


\par\indent Ich habe diese Beobachtungen mit meinen letzten parabolischen Elementen verglichen, so wie Herr \textsc{Nicolai} die letzten \textsc{Olbers}schen und Seeberger Beobachtungen. Folgendes sind die Resultate:

\Needspace{20\baselineskip}\begin{center}Bremer Beobachtungen\end{center}
\begin{center}\fontsize{10}{12}\selectfont\setlength{\tabcolsep}{8pt}\renewcommand{\arraystretch}{1.15}\begin{tabular}{l|c|c}
& \multicolumn{2}{c}{Unterschied} \\
Tag der & & \\
Beobacht. & in AR. & in Decl. \\ \hline
Octob. 19 & $-73''$ 
 & $+152''$ \\
19 & $-73$ & ---\;--- \\
20 & $-39$ & $+139$ \\
24 & $-82$ & $+227$ \\
24 & $-84$ & $+205$ \\
25 & $-108$ & ---\;--- \\
25 & $-91$ & ---\;--- \\
28 & $-60$ & $+293$ \\
Nov. 4 & $-159$ & ---\;--- \\
4 & $-155$ & $+342$ \\
7 & $-91$ & ---\;--- \\
7 & ---\;--- & $+377$ \\
7 & $-102$ & ---\;--- \\
9 & $-274$ & $+379$ \\
\end{tabular}\end{center}
\Needspace{15\baselineskip}\begin{center}Seeberger Beobachtungen\end{center}
\begin{center}\fontsize{10}{12}\selectfont\setlength{\tabcolsep}{8pt}\renewcommand{\arraystretch}{1.15}\begin{tabular}{l|c|c}
& \multicolumn{2}{c}{Unterschied} \\
Tag der & & \\
Beobacht. & in AR. & in Decl. \\ \hline
Octob. 20 & $-113''$ & $+196''$ \\
24 & $-109$ & $+175$ \\
25 & $-88$ & $+263$ \\
28 & $-87$ & $+327$ \\
29 & $-124$ & $+265$ \\
Nov. 4 & $-110$ & $+387$ \\
5 & $-74$ & ---\;--- \\
\end{tabular}\end{center}


\Needspace{19\baselineskip}\par\indent Doctor \textsc{Olbers} hat mir die ersten Beobachtungen des Herrn \textsc{Flaugergues} mitgetheilt, denen ich Herrn \textsc{Nicolai}s Berechnung der Unterschiede von den Elementen beifüge

\Needspace{15\baselineskip}
\begin{center}\fontsize{9}{11}\selectfont\setlength{\tabcolsep}{3pt}\renewcommand{\arraystretch}{1.15}\begin{tabular}{l|c|c|c|c}
\shortstack{Tag der\\Beobacht.} & \shortstack{AR.\\Cometae} & \shortstack{Abweichung\\v. d. Element.} & \shortstack{Declinatio\\austr. Comet.} & \shortstack{Abweichung\\v. d. Element.}\\\hline
März 26 & $120^\circ\,16'\,0''$ & $+859''$ & $29^\circ\,15'\,0''$ & $+300''$\\
28 & $119\;52\;56$ & $+316$ & $28\;\;7\;\;0$ & $+257$\\
29 & $119\;41\;\;4$ & $+127$ & $27\;32\;57$ & $+226$\\
30 & $119\;29\;26$ & $-50$ & $26\;58\;22$ & $+193$\\
31 & $119\;18\;36$ & $-306$ & $26\;23\;13$ & $+297$\\
April 1 & $119\;\;7\;38$ & $-382$ & $25\;50\;17$ & $+219$\\
\end{tabular}\end{center}

\SourceLeaf{347}\par\indent Herrn \textsc{Flaugergues} erste Beobachtung war in frühern Nachrichten in AR. um $10'$ grösser angegeben, und dann wäre der Unterschied nur $259''$. Ich enthalte mich noch über die grossen Differenzen zu urtheilen, bis ich einmal Zeit gewinne, selbst wieder über den Cometen Rechnungen anzustellen, wozu ich jetzt noch nicht kommen kann.


